% =====================================================================
%  CHƯƠNG 2. GIẢI THÍCH CƠ SỞ LÝ THUYẾT CỦA G, H VÀ MỞ RỘNG SANG MÃ (8,4)
% =====================================================================
\chapter{GIẢI THÍCH CƠ SỞ LÝ THUYẾT CỦA G, H VÀ MỞ RỘNG SANG MÃ (8,4)}
\label{chap:2}

Trong chương này, $h_j$ (cột thứ $j$ của $H$) được viết như một vectơ hàng
ba thành phần. Với $x = (x_1\,\dots\,x_7) \in \GFn{7}$, tích $x\HT$ là tổng
các hàng của $\HT$ ứng với các vị trí $x_j = 1$:
\begin{equation}\label{eq:tongcot}
  x\HT = \bigoplus_{j\,:\,x_j = 1} h_j .
\end{equation}

% ---------------------------------------------------------------------
\section{Vì sao \texorpdfstring{$G$}{G} có dạng \texorpdfstring{$[I_4 \mid P]$}{[I4 | P]}}
\label{sec:2.1}

Bốn cột đầu của $G$ tạo thành $I_4$, nên bốn hàng của $G$ độc lập tuyến tính
và $\rank G = 4$. Do đó ánh xạ $u \mapsto uG$ từ $\GFn{4}$ vào $\GFn{7}$ là
đơn ánh, và $C$ (không gian hàng của $G$) có số chiều $4$, gồm đúng
$2^4 = 16$ từ mã phân biệt. Hơn nữa $uG = (u \mid uP)$, nên dữ liệu đọc
thẳng từ 4 bit đầu của từ mã.

\canviet{Giải thích thêm: mọi mã tuyến tính đều đưa được về dạng hệ thống
bằng các phép biến đổi sơ cấp trên hàng (và hoán vị cột nếu cần); vai trò của
$P$ (mỗi cột của $P$ cho một phương trình bit kiểm tra). Liên hệ hạng ma
trận, không gian hàng.}

% ---------------------------------------------------------------------
\section{Vì sao \texorpdfstring{$H = [P^{\mathsf T} \mid I_3]$ và $GH^{\mathsf T} = 0$}{H = [PT | I3] và GHT = 0}}
\label{sec:2.2}

\begin{menhde}\label{md:GHT}
$G\HT = 0$ và $C = \{\, c \in \GFn{7} : c\HT = 0 \,\}$.
\end{menhde}
\begin{proof}
Nhân theo khối, với $\HT = \begin{bmatrix} P \\ I_3 \end{bmatrix}$:
\[
  G\HT = [\,I_4 \mid P\,]\begin{bmatrix} P \\ I_3 \end{bmatrix}
       = I_4 P \xor P I_3 = P \xor P = 0 ,
\]
vì trên $\GF$ mọi phần tử cộng với chính nó bằng $0$. Do đó với mọi từ mã
$c = uG$ ta có $c\HT = u(G\HT) = 0$, tức $C \subseteq \ker$, với
$\ker = \{x \in \GFn{7} : x\HT = 0\}$.

Ba cột cuối của $H$ tạo thành $I_3$ nên $\rank H = 3$. Theo định lý hạng --
số chiều hạt nhân, $\dim \ker = 7 - 3 = 4 = \dim C$. Không gian con $C$ nằm
trong $\ker$ và cùng số chiều nên $C = \ker$.
\end{proof}

\canviet{Giải thích vì sao chọn $H = [\PT \mid I_3]$: đây là cách viết ba
phương trình bit kiểm tra $uP \xor (p_1\,p_2\,p_3) = 0$ thành ma trận. Lưu ý
dấu: trên $\GF$ thì $-\PT = \PT$.}

% ---------------------------------------------------------------------
\section{Vì sao các cột của \texorpdfstring{$H$}{H} khác 0 và đôi một khác nhau}
\label{sec:2.3}

\canviet{(1) Nếu có $h_j = 0$ thì lỗi ở bit $j$ cho $s = 0$, không bị phát
hiện. (2) Nếu có $h_i = h_j$ với $i \neq j$ thì lỗi ở bit $i$ và lỗi ở bit
$j$ cho cùng hội chứng, không phân biệt được. (3) $\GFn{3}$ có đúng $2^3 - 1 = 7$
vectơ khác $0$, nên $n = 7$ là độ dài lớn nhất với 3 bit kiểm tra; tổng quát
mã Hamming có $n = 2^m - 1$, $k = 2^m - 1 - m$.}

% ---------------------------------------------------------------------
\section{\texorpdfstring{$d_{\min} = 3$}{dmin = 3} suy trực tiếp từ \texorpdfstring{$H$}{H}}
\label{sec:2.4}

\begin{dinhly}\label{dl:dmin3}
Mã Hamming (7,4) có $\dmin = 3$.
\end{dinhly}
\begin{proof}
Theo Mệnh đề~\ref{md:GHT} và \eqref{eq:tongcot}, $c \in C$ khi và chỉ khi
tổng các cột $h_j$ ứng với các vị trí $c_j = 1$ bằng $0$.
\begin{itemize}
  \item Không có từ mã trọng số $1$: điều đó đòi hỏi $h_j = 0$, trái với việc
        mọi cột của $H$ khác $0$.
  \item Không có từ mã trọng số $2$: điều đó đòi hỏi $h_i \xor h_j = 0$, tức
        $h_i = h_j$ với $i \neq j$, trái với việc các cột đôi một khác nhau.
  \item Có từ mã trọng số $3$: $h_3 \xor h_6 \xor h_7
        = \vt{0,1,1} \xor \vt{0,1,0} \xor \vt{0,0,1} = \vt{0,0,0}$, nên
        $c = (0\,0\,1\,0\,0\,1\,1) \in C$ (hàng 3 của Bảng~\ref{tab:16tuma}).
\end{itemize}
Theo Mệnh đề~\ref{md:dmin-trongso}, $\dmin = 3$.
\end{proof}

\canviet{Hệ quả theo Định lý~\ref{dl:phathien-sua}: sửa được $t = 1$ lỗi,
hoặc (nếu chỉ dùng để phát hiện) phát hiện được tối đa $2$ lỗi. Tổng quát:
$\dmin$ là số nhỏ nhất các cột phụ thuộc tuyến tính của $H$.}

% ---------------------------------------------------------------------
% LƯU Ý TÍCH HỢP: tiêu đề giữ nguyên theo khung BTLDSTT.pdf. Tài liệu
% "quy_uoc_chung_dang_he_thong" mục 8 đề xuất đổi thành
% "Mã hoàn hảo: sửa được 1 lỗi, nhưng không thể đồng thời sửa 1 lỗi và
%  phát hiện 2 lỗi" — CHỜ NHÓM TRƯỞNG XÁC NHẬN.
\section{Mã hoàn hảo: vì sao sửa được 1 lỗi nhưng không phát hiện được 2 lỗi}
\label{sec:2.5}

Với mỗi từ mã $c$, gọi $B(c) = \{x \in \GFn{7} : \dist{x}{c} \le 1\}$ là
quả cầu Hamming bán kính $1$ tâm $c$; $B(c)$ gồm $c$ và $7$ vectơ sai khác
$c$ đúng một bit, nên $|B(c)| = 1 + 7 = 8$.

\begin{dinhly}[Mã hoàn hảo]\label{dl:hoanhao}
Các quả cầu $B(c)$, $c \in C$, đôi một rời nhau và phủ kín $\GFn{7}$.
\end{dinhly}
\begin{proof}
Nếu $x \in B(c) \cap B(c')$ với $c \neq c'$ thì
$\dist{c}{c'} \le \dist{c}{x} + \dist{x}{c'} \le 2 < \dmin = 3$, vô lý.
Vậy các quả cầu rời nhau, và hợp của chúng có
$16 \cdot 8 = 128 = 2^7 = |\GFn{7}|$ phần tử, tức phủ kín $\GFn{7}$.
\end{proof}

Hệ quả: mọi vectơ nhận được $r$ cách đúng một từ mã không quá một bit, nên
bộ giải mã luôn trả về một từ mã. Nếu có đúng hai lỗi, $r = c \xor e$ với
$\wt{e} = 2$ thì $\dist{r}{c} = 2$, nên $r$ nằm trong quả cầu của một từ mã
$c' \neq c$ và bộ giải mã chọn $c'$: \emph{sửa nhầm}. Về hội chứng,
$s = h_i \xor h_j \neq 0$ (vì $h_i \neq h_j$) và là một vectơ khác $0$ của
$\GFn{3}$, nên trùng với một cột $h_k$ nào đó, giống hệt hội chứng của một
lỗi đơn ở bit $k$.

\begin{nhanxet}
Không đánh đồng khả năng sửa một lỗi với khả năng phát hiện hai lỗi. Nếu
\emph{chỉ phát hiện} (không sửa), mọi mẫu hai lỗi đều cho $s \neq 0$ và bị
phát hiện, đúng với Định lý~\ref{dl:phathien-sua}(a). Nhưng khi bộ giải mã
\emph{sửa} theo Bảng~\ref{tab:hoichung}, nó không phân biệt được một lỗi với
hai lỗi, nên không thể đồng thời sửa một lỗi và phát hiện hai lỗi. Muốn
làm được cả hai cần mã mở rộng (8,4) ở Mục~\ref{sec:2.7}.
\end{nhanxet}

\canviet{Giới thiệu cận Hamming tổng quát $2^k \sum_{i=0}^{t}\binom{n}{i} \le 2^n$
và dấu bằng với $(n,k,t) = (7,4,1)$.}

% ---------------------------------------------------------------------
\section{Dạng cổ điển: hội chứng chính là số thứ tự vị trí lỗi}
\label{sec:2.6}

Trong dạng cổ điển, bit kiểm tra đặt ở các vị trí lũy thừa của 2:
$c_{\text{cổ điển}} = (p_1\,p_2\,u_1\,p_3\,u_2\,u_3\,u_4)$, và cột thứ $j$ của
ma trận kiểm tra là số $j$ viết ở hệ nhị phân (hàng 1 là bit thấp nhất):
\begin{equation}
  H_{\text{cổ điển}} = \MatHcd .
\end{equation}
Khi có một lỗi ở vị trí $j$, hội chứng $s = (s_1\,s_2\,s_3)$ là cột thứ $j$,
nên đọc thẳng ra vị trí lỗi:
\begin{equation}
  j = s_1 + 2s_2 + 4s_3 .
\end{equation}

Dạng cổ điển và dạng hệ thống là cùng một mã sau một hoán vị vị trí: nếu
$c = (c_1,\dots,c_7)$ là từ mã cổ điển thì
\begin{equation}\label{eq:hoanvi}
  c_{\text{hệ thống}} = (c_3,\,c_5,\,c_6,\,c_7,\,c_1,\,c_2,\,c_4).
\end{equation}
Ví dụ $0110011$ (cổ điển) $\leftrightarrow$ $1011010$ (hệ thống). Hội
chứng của cùng một bit là cùng một vectơ ở hai dạng; chỉ khác nhãn vị trí
(Bảng~\ref{tab:codien}). Báo cáo và phần mềm dùng \emph{dạng hệ thống}, nên
hội chứng được tra theo Bảng~\ref{tab:hoichung}.

\begin{table}[H]
\centering
\caption{Đối chiếu vị trí giữa dạng cổ điển và dạng hệ thống}
\label{tab:codien}
\begin{tabular}{|c|c|c|c|}
\hline
\textbf{Vị trí cổ điển $j$} & \textbf{Bit} & \textbf{Hội chứng $(s_1 s_2 s_3)$} &
\textbf{Vị trí hệ thống}\\ \hline
1 & $p_1$ & 100 & 5\\
2 & $p_2$ & 010 & 6\\
3 & $u_1$ & 110 & 1\\
4 & $p_3$ & 001 & 7\\
5 & $u_2$ & 101 & 2\\
6 & $u_3$ & 011 & 3\\
7 & $u_4$ & 111 & 4\\ \hline
\end{tabular}
\end{table}

\canviet{Giải thích bằng lời ưu điểm của mỗi dạng: dạng cổ điển đọc vị trí
lỗi trực tiếp; dạng hệ thống tách dữ liệu dễ (4 bit đầu) và gắn với
$G = [I_4 \mid P]$.}

% ---------------------------------------------------------------------
\section{Mã Hamming mở rộng (8,4) để phát hiện hai lỗi}
\label{sec:2.7}

Thêm vào mỗi từ mã (7,4) một \emph{bit kiểm tra tổng thể}
\begin{equation}
  c_8 = c_1 \xor c_2 \xor \dots \xor c_7 ,
\end{equation}
ta được từ mã mở rộng $c' = (c_1\,\dots\,c_7\,c_8)$ có trọng số chẵn. Ma trận
kiểm tra mở rộng cỡ $4 \times 8$ gồm ba hàng của $H$ thêm cột $0$ ở cuối và
một hàng toàn số $1$:
\begin{equation}\label{eq:Hmr}
  \Hmr = \MatHmr .
\end{equation}

Khi nhận $r' = (r_1\,\dots\,r_8)$, tính
\begin{equation}
  s = (r_1\,\dots\,r_7)\HT, \qquad q = r_1 \xor r_2 \xor \dots \xor r_8 ,
\end{equation}
rồi kết luận theo Bảng~\ref{tab:84}.

\begin{table}[H]
\centering
\caption{Quy tắc giải mã của mã Hamming mở rộng (8,4)}
\label{tab:84}
\begin{tabular}{|c|c|l|}
\hline
$s$ & $q$ & \textbf{Kết luận}\\ \hline
$= 000$   & 0 & Không phát hiện lỗi\\
$\ne 000$ & 1 & Một lỗi trong 7 bit đầu: lật bit $j$ với $s = h_j$\\
$= 000$   & 1 & Lỗi ở bit thứ 8: lật bit 8\\
$\ne 000$ & 0 & Phát hiện lỗi kép, không sửa\\ \hline
\end{tabular}
\end{table}

Lý do: một lỗi làm $q = 1$ (số bit 1 đổi tính chẵn lẻ), và $s = h_j \ne 0$
nếu lỗi nằm trong 7 bit đầu, $s = 0$ nếu lỗi ở bit 8. Hai lỗi giữ $q = 0$;
nếu cả hai nằm trong 7 bit đầu thì $s = h_i \xor h_j \ne 0$, nếu một lỗi ở
bit 8 thì $s = h_i \ne 0$. Vậy hai lỗi luôn rơi vào dòng cuối của bảng, không
bị nhầm với một lỗi.

\canviet{Chứng minh mã mở rộng có $\dmin = 4$ (mọi từ mã có trọng số chẵn,
nên các từ mã trọng số 3 thành trọng số 4) và suy ra: sửa 1 lỗi đồng thời phát
hiện 2 lỗi. Nêu giá phải trả: tỉ lệ mã giảm từ $4/7$ xuống $4/8 = 1/2$.}
